% Agentic Coding for Economists: EIEF / Luiss eight-hour workshop
% Fourteen ideas, one example (synthetic Card--Krueger panel).
% Build: latexmk -pdf -interaction=nonstopmode workshop_slides.tex
% (from slides/workshop/)

% Shared Beamer preamble — Agentic Coding for Economists (EIEF / Luiss)
% Include via: % Shared Beamer preamble — Agentic Coding for Economists (EIEF / Luiss)
% Include via: % Shared Beamer preamble — Agentic Coding for Economists (EIEF / Luiss)
% Include via: \input{../shared/beamer_preamble}

\documentclass[aspectratio=169,11pt]{beamer}

\usetheme{Madrid}
\usecolortheme{default}

\definecolor{courseblue}{RGB}{0,91,150}
\definecolor{courseteal}{RGB}{0,150,136}
\definecolor{coursegray}{RGB}{100,100,100}
\definecolor{coursebg}{RGB}{245,245,250}
\definecolor{codegreen}{RGB}{40,160,80}
\definecolor{codepurple}{RGB}{120,60,180}
\definecolor{codered}{RGB}{200,50,50}

\setbeamercolor{structure}{fg=courseblue}
\setbeamercolor{palette primary}{bg=courseblue,fg=white}
\setbeamercolor{palette secondary}{bg=courseteal,fg=white}
\setbeamercolor{title}{fg=courseblue}
\setbeamercolor{frametitle}{fg=courseblue,bg=white}
\setbeamercolor{block title}{bg=courseblue,fg=white}
\setbeamercolor{block body}{bg=coursebg}
\setbeamercolor{item}{parent=normal text}
\setbeamercolor{itemize item}{fg=courseblue}
\setbeamercolor{itemize subitem}{fg=courseblue!85!black}
\setbeamercolor{itemize subsubitem}{fg=courseblue!70!black}
\setbeamercolor{enumerate item}{fg=courseblue}
\setbeamercolor{enumerate subitem}{fg=courseblue!85!black}
\setbeamercolor{itemize/enumerate body}{parent=normal text}
\setbeamercolor{itemize/enumerate subbody}{parent=normal text}
\setbeamercolor{itemize/enumerate subsubbody}{parent=normal text}
\setbeamercolor{alerted text}{fg=codered}

\usefonttheme{professionalfonts}
\setbeamerfont{title}{size=\Large,series=\bfseries}
\setbeamerfont{frametitle}{size=\large,series=\bfseries}

\usepackage[utf8]{inputenc}
\usepackage[T1]{fontenc}
\usepackage{textcomp}
\usepackage{amsmath,amssymb}
\usepackage{booktabs}
\usepackage{graphicx}
\usepackage{hyperref}
\usepackage{listings}
\usepackage{tikz}
\usetikzlibrary{arrows.meta,positioning,shapes.geometric,calc,fit}

\hypersetup{colorlinks=true,linkcolor=courseblue,urlcolor=courseblue,citecolor=courseblue}

\lstset{
  basicstyle=\ttfamily\scriptsize,
  keywordstyle=\color{courseblue}\bfseries,
  commentstyle=\color{codegreen}\itshape,
  stringstyle=\color{codepurple},
  breaklines=true,
  frame=single,
  backgroundcolor=\color{coursebg},
  rulecolor=\color{coursegray!30},
  numbers=none,
  showstringspaces=false,
  tabsize=2,
}

\lstdefinestyle{python}{language=Python,morekeywords={self,True,False,None}}
\lstdefinestyle{bash}{language=bash,morekeywords={git,codex,cursor,python,pip,claude,latexmk,pytest}}
\lstdefinestyle{markdown}{basicstyle=\ttfamily\scriptsize,frame=single,backgroundcolor=\color{coursebg}}

\newcommand{\concept}[1]{\textcolor{courseblue}{\textbf{#1}}}
\newcommand{\artifact}[1]{\texttt{\textcolor{courseteal}{#1}}}
\newcommand{\timebox}[1]{\textcolor{coursegray}{\small\textsf{[#1]}}}
\newcommand{\prompttext}[1]{%
  \begin{minipage}{\linewidth}%
    \colorbox{coursebg}{\parbox{\dimexpr\linewidth-2\fboxsep}{\ttfamily\footnotesize #1}}%
  \end{minipage}%
}
\newcommand{\lanetag}[1]{\fcolorbox{courseblue}{white}{\small\textsf{#1}}}

% Public participant pack. The title slide asks the agent to fork, then confirm a folder before cloning the fork.
\newcommand{\WorkshopRepoUrl}{https://github.com/meleantonio/ac4e-eief-workshop}

\setbeamertemplate{navigation symbols}{}
\setbeamertemplate{footline}{%
  \leavevmode%
  \hbox{%
    \begin{beamercolorbox}[wd=.33\paperwidth,ht=2.25ex,dp=1ex,center]{author in head/foot}%
      \usebeamerfont{author in head/foot}\insertshortauthor
    \end{beamercolorbox}%
    \begin{beamercolorbox}[wd=.34\paperwidth,ht=2.25ex,dp=1ex,center]{title in head/foot}%
      \usebeamerfont{title in head/foot}\insertshorttitle
    \end{beamercolorbox}%
    \begin{beamercolorbox}[wd=.33\paperwidth,ht=2.25ex,dp=1ex,right]{date in head/foot}%
      \usebeamerfont{date in head/foot}\insertshortdate{}\hspace*{2em}
      \insertframenumber{}/\inserttotalframenumber\hspace*{2ex}
    \end{beamercolorbox}}%
  \vskip0pt%
}

\author{Antonio Mele}
\institute{EIEF / Luiss}
\date{26--27 November 2026}


\documentclass[aspectratio=169,11pt]{beamer}

\usetheme{Madrid}
\usecolortheme{default}

\definecolor{courseblue}{RGB}{0,91,150}
\definecolor{courseteal}{RGB}{0,150,136}
\definecolor{coursegray}{RGB}{100,100,100}
\definecolor{coursebg}{RGB}{245,245,250}
\definecolor{codegreen}{RGB}{40,160,80}
\definecolor{codepurple}{RGB}{120,60,180}
\definecolor{codered}{RGB}{200,50,50}

\setbeamercolor{structure}{fg=courseblue}
\setbeamercolor{palette primary}{bg=courseblue,fg=white}
\setbeamercolor{palette secondary}{bg=courseteal,fg=white}
\setbeamercolor{title}{fg=courseblue}
\setbeamercolor{frametitle}{fg=courseblue,bg=white}
\setbeamercolor{block title}{bg=courseblue,fg=white}
\setbeamercolor{block body}{bg=coursebg}
\setbeamercolor{item}{parent=normal text}
\setbeamercolor{itemize item}{fg=courseblue}
\setbeamercolor{itemize subitem}{fg=courseblue!85!black}
\setbeamercolor{itemize subsubitem}{fg=courseblue!70!black}
\setbeamercolor{enumerate item}{fg=courseblue}
\setbeamercolor{enumerate subitem}{fg=courseblue!85!black}
\setbeamercolor{itemize/enumerate body}{parent=normal text}
\setbeamercolor{itemize/enumerate subbody}{parent=normal text}
\setbeamercolor{itemize/enumerate subsubbody}{parent=normal text}
\setbeamercolor{alerted text}{fg=codered}

\usefonttheme{professionalfonts}
\setbeamerfont{title}{size=\Large,series=\bfseries}
\setbeamerfont{frametitle}{size=\large,series=\bfseries}

\usepackage[utf8]{inputenc}
\usepackage[T1]{fontenc}
\usepackage{textcomp}
\usepackage{amsmath,amssymb}
\usepackage{booktabs}
\usepackage{graphicx}
\usepackage{hyperref}
\usepackage{listings}
\usepackage{tikz}
\usetikzlibrary{arrows.meta,positioning,shapes.geometric,calc,fit}

\hypersetup{colorlinks=true,linkcolor=courseblue,urlcolor=courseblue,citecolor=courseblue}

\lstset{
  basicstyle=\ttfamily\scriptsize,
  keywordstyle=\color{courseblue}\bfseries,
  commentstyle=\color{codegreen}\itshape,
  stringstyle=\color{codepurple},
  breaklines=true,
  frame=single,
  backgroundcolor=\color{coursebg},
  rulecolor=\color{coursegray!30},
  numbers=none,
  showstringspaces=false,
  tabsize=2,
}

\lstdefinestyle{python}{language=Python,morekeywords={self,True,False,None}}
\lstdefinestyle{bash}{language=bash,morekeywords={git,codex,cursor,python,pip,claude,latexmk,pytest}}
\lstdefinestyle{markdown}{basicstyle=\ttfamily\scriptsize,frame=single,backgroundcolor=\color{coursebg}}

\newcommand{\concept}[1]{\textcolor{courseblue}{\textbf{#1}}}
\newcommand{\artifact}[1]{\texttt{\textcolor{courseteal}{#1}}}
\newcommand{\timebox}[1]{\textcolor{coursegray}{\small\textsf{[#1]}}}
\newcommand{\prompttext}[1]{%
  \begin{minipage}{\linewidth}%
    \colorbox{coursebg}{\parbox{\dimexpr\linewidth-2\fboxsep}{\ttfamily\footnotesize #1}}%
  \end{minipage}%
}
\newcommand{\lanetag}[1]{\fcolorbox{courseblue}{white}{\small\textsf{#1}}}

% Public participant pack. The title slide asks the agent to fork, then confirm a folder before cloning the fork.
\newcommand{\WorkshopRepoUrl}{https://github.com/meleantonio/ac4e-eief-workshop}

\setbeamertemplate{navigation symbols}{}
\setbeamertemplate{footline}{%
  \leavevmode%
  \hbox{%
    \begin{beamercolorbox}[wd=.33\paperwidth,ht=2.25ex,dp=1ex,center]{author in head/foot}%
      \usebeamerfont{author in head/foot}\insertshortauthor
    \end{beamercolorbox}%
    \begin{beamercolorbox}[wd=.34\paperwidth,ht=2.25ex,dp=1ex,center]{title in head/foot}%
      \usebeamerfont{title in head/foot}\insertshorttitle
    \end{beamercolorbox}%
    \begin{beamercolorbox}[wd=.33\paperwidth,ht=2.25ex,dp=1ex,right]{date in head/foot}%
      \usebeamerfont{date in head/foot}\insertshortdate{}\hspace*{2em}
      \insertframenumber{}/\inserttotalframenumber\hspace*{2ex}
    \end{beamercolorbox}}%
  \vskip0pt%
}

\author{Antonio Mele}
\institute{EIEF / Luiss}
\date{26--27 November 2026}


\documentclass[aspectratio=169,11pt]{beamer}

\usetheme{Madrid}
\usecolortheme{default}

\definecolor{courseblue}{RGB}{0,91,150}
\definecolor{courseteal}{RGB}{0,150,136}
\definecolor{coursegray}{RGB}{100,100,100}
\definecolor{coursebg}{RGB}{245,245,250}
\definecolor{codegreen}{RGB}{40,160,80}
\definecolor{codepurple}{RGB}{120,60,180}
\definecolor{codered}{RGB}{200,50,50}

\setbeamercolor{structure}{fg=courseblue}
\setbeamercolor{palette primary}{bg=courseblue,fg=white}
\setbeamercolor{palette secondary}{bg=courseteal,fg=white}
\setbeamercolor{title}{fg=courseblue}
\setbeamercolor{frametitle}{fg=courseblue,bg=white}
\setbeamercolor{block title}{bg=courseblue,fg=white}
\setbeamercolor{block body}{bg=coursebg}
\setbeamercolor{item}{parent=normal text}
\setbeamercolor{itemize item}{fg=courseblue}
\setbeamercolor{itemize subitem}{fg=courseblue!85!black}
\setbeamercolor{itemize subsubitem}{fg=courseblue!70!black}
\setbeamercolor{enumerate item}{fg=courseblue}
\setbeamercolor{enumerate subitem}{fg=courseblue!85!black}
\setbeamercolor{itemize/enumerate body}{parent=normal text}
\setbeamercolor{itemize/enumerate subbody}{parent=normal text}
\setbeamercolor{itemize/enumerate subsubbody}{parent=normal text}
\setbeamercolor{alerted text}{fg=codered}

\usefonttheme{professionalfonts}
\setbeamerfont{title}{size=\Large,series=\bfseries}
\setbeamerfont{frametitle}{size=\large,series=\bfseries}

\usepackage[utf8]{inputenc}
\usepackage[T1]{fontenc}
\usepackage{textcomp}
\usepackage{amsmath,amssymb}
\usepackage{booktabs}
\usepackage{graphicx}
\usepackage{hyperref}
\usepackage{listings}
\usepackage{tikz}
\usetikzlibrary{arrows.meta,positioning,shapes.geometric,calc,fit}

\hypersetup{colorlinks=true,linkcolor=courseblue,urlcolor=courseblue,citecolor=courseblue}

\lstset{
  basicstyle=\ttfamily\scriptsize,
  keywordstyle=\color{courseblue}\bfseries,
  commentstyle=\color{codegreen}\itshape,
  stringstyle=\color{codepurple},
  breaklines=true,
  frame=single,
  backgroundcolor=\color{coursebg},
  rulecolor=\color{coursegray!30},
  numbers=none,
  showstringspaces=false,
  tabsize=2,
}

\lstdefinestyle{python}{language=Python,morekeywords={self,True,False,None}}
\lstdefinestyle{bash}{language=bash,morekeywords={git,codex,cursor,python,pip,claude,latexmk,pytest}}
\lstdefinestyle{markdown}{basicstyle=\ttfamily\scriptsize,frame=single,backgroundcolor=\color{coursebg}}

\newcommand{\concept}[1]{\textcolor{courseblue}{\textbf{#1}}}
\newcommand{\artifact}[1]{\texttt{\textcolor{courseteal}{#1}}}
\newcommand{\timebox}[1]{\textcolor{coursegray}{\small\textsf{[#1]}}}
\newcommand{\prompttext}[1]{%
  \begin{minipage}{\linewidth}%
    \colorbox{coursebg}{\parbox{\dimexpr\linewidth-2\fboxsep}{\ttfamily\footnotesize #1}}%
  \end{minipage}%
}
\newcommand{\lanetag}[1]{\fcolorbox{courseblue}{white}{\small\textsf{#1}}}

% Public participant pack. The title slide asks the agent to fork, then confirm a folder before cloning the fork.
\newcommand{\WorkshopRepoUrl}{https://github.com/meleantonio/ac4e-eief-workshop}

\setbeamertemplate{navigation symbols}{}
\setbeamertemplate{footline}{%
  \leavevmode%
  \hbox{%
    \begin{beamercolorbox}[wd=.33\paperwidth,ht=2.25ex,dp=1ex,center]{author in head/foot}%
      \usebeamerfont{author in head/foot}\insertshortauthor
    \end{beamercolorbox}%
    \begin{beamercolorbox}[wd=.34\paperwidth,ht=2.25ex,dp=1ex,center]{title in head/foot}%
      \usebeamerfont{title in head/foot}\insertshorttitle
    \end{beamercolorbox}%
    \begin{beamercolorbox}[wd=.33\paperwidth,ht=2.25ex,dp=1ex,right]{date in head/foot}%
      \usebeamerfont{date in head/foot}\insertshortdate{}\hspace*{2em}
      \insertframenumber{}/\inserttotalframenumber\hspace*{2ex}
    \end{beamercolorbox}}%
  \vskip0pt%
}

\author{Antonio Mele}
\institute{EIEF / Luiss}
\date{26--27 November 2026}


\title{Agentic Coding for Economists}
\subtitle{EIEF / Luiss: eight hours, fourteen ideas}
\author{Antonio Mele}
\institute{EIEF / Luiss}
\date{26--27 November 2026}

% Concept frames: definition block, then the everyday comparison.
\newcommand{\inplainwords}[1]{%
  \begin{block}{In plain words}\large #1\end{block}}
\newcommand{\likethis}[1]{%
  \vspace{0.8em}{\large \textcolor{courseteal}{\textbf{Think of it as:}} #1}}
\newcommand{\youtry}[1]{\textcolor{coursegray}{\small\textsf{Your turn \textperiodcentered\ #1}}}

\begin{document}

% =============================================================================
\section{Welcome}
% =============================================================================

\begin{frame}{}
  \begin{center}
    {\LARGE \textbf{Agentic Coding for Economists}}\\[0.12em]
    {\large EIEF / Luiss: eight hours, fourteen ideas}\\[0.25em]
    {\small Thu 26 Nov 2026 \quad 09:00--12:00 \quad\textperiodcentered\quad
      Fri 27 Nov 2026 \quad 09:30--12:00 and 14:00--16:30 (Rome)}\\[0.18em]
    {\small Antonio Mele \quad\textperiodcentered\quad
      Claude Code \textperiodcentered\ Codex and Cursor work the same way}\\[0.1em]
    {\small One example all day: the synthetic Card--Krueger panel}\\[0.35em]
    {\footnotesize Paste this into Claude Code, Codex, or Cursor.}
  \end{center}
  \vspace{0.25em}
  \begin{center}
    \fcolorbox{coursegray!35}{coursebg}{%
      \begin{minipage}{0.94\linewidth}
        \ttfamily\scriptsize
        \mbox{Go to \WorkshopSiteUrl{} and follow instructions.}
      \end{minipage}}
  \end{center}
\end{frame}

\begin{frame}{Fourteen ideas on the Rome clock}
  \small
  \begin{tabular}{@{}lll@{}}
    \toprule
    \textbf{Slot} & \textbf{Clock (Rome)} & \textbf{Ideas} \\
    \midrule
    1 & Thu 09:00--11:00 & Agent vs chatbot \textperiodcentered\ Version control \\
      & & \artifact{AGENTS.md}: boundaries first \textperiodcentered\ To-do lists \\
    \addlinespace[0.3em]
    2 & Thu 11:00--12:00 & Skills: why \textperiodcentered\ Create a Skill \\
      & Fri 09:30--10:15 & Sub-agent vs Skill \textperiodcentered\ Skill or sub-agent? \\
    \addlinespace[0.3em]
    3 & Fri 10:15--12:00 & Using sub-agents \textperiodcentered\ Orchestrating and swarms \\
      & & Issue-board swarm (agents work over lunch) \\
    \addlinespace[0.3em]
    4 & Fri 14:00--16:30 & Goals \textperiodcentered\ Acceptance criteria \\
      & & Verification \textperiodcentered\ Capstone \textperiodcentered\ 7-day plan \\
    \bottomrule
  \end{tabular}
  \vspace{0.5em}

  {\small Each idea: definition, everyday comparison, Card--Krueger example, your turn.}
\end{frame}

\begin{frame}[fragile]{One example all day: Card--Krueger}
  \begin{itemize}
    \item April 1992: New Jersey raises its minimum wage; Pennsylvania does not.
    \item Compare fast-food jobs before and after: the NJ--PA difference-in-differences.
    \item Our folder: \artifact{examples/card-krueger-toy/}.
    \item \textbf{Synthetic teaching data}, not the real survey. No research claims.
  \end{itemize}
  \vspace{0.3em}
  The check command we run all day:
  \begin{lstlisting}[style=bash]
python3 -m pytest examples/card-krueger-toy/tests examples/mini-economics/tests
  \end{lstlisting}
\end{frame}

\begin{frame}[fragile]{Room rules}
  \timebox{09:00--09:10}
  \begin{itemize}
    \item You are the manager; the agent is a fast RA.
    \item Every agent change is a \textbf{proposal} until you read it and the check passes.
    \item Agents never certify identification or causality. That is your job.
    \item No confidential data, referee files or passwords in any prompt.
  \end{itemize}
  \vspace{0.3em}
  First, check your laptop:
  \begin{lstlisting}[style=bash]
python3 scripts/validate_setup.py
  \end{lstlisting}
\end{frame}

% Slot 1: C1--C4 (Thu 26 Nov, 09:00--11:00 Rome)

\section{Slot 1: C1--C4}

\begin{frame}{Slot 1 \textperiodcentered\ Thu 09:00--11:00}
  \begin{center}
    \begin{tabular}{@{}rl@{}}
      \toprule
      \textbf{Clock} & \textbf{Block} \\
      \midrule
      09:00--09:10 & Welcome, room rules, laptop check \\
      09:10--09:30 & \concept{C1} Agent vs chatbot \\
      09:30--09:55 & \concept{C2} Version control: your first pull request \\
      09:55--10:35 & \concept{C3} \artifact{AGENTS.md}: boundaries first, with privacy \\
      10:35--10:45 & Pause / set-up help \\
      10:45--11:00 & \concept{C4} To-do lists \\
      \bottomrule
    \end{tabular}
  \end{center}
\end{frame}

% ---------------------------------------------------------------- C1
\begin{frame}{C1 \textperiodcentered\ Agent vs chatbot}
  \timebox{09:10--09:30}
  \inplainwords{A chatbot only talks. An agent can also act: it opens your files,
    runs your code, reads the result and tries again.}
  \likethis{phoning a friend for advice vs hiring an RA who sits at your desk.}
\end{frame}

\begin{frame}{C1 \textperiodcentered\ Card--Krueger: same question, two answers}
  \centering
  {\large ``What is the NJ--PA change in employment?''}
  \vspace{0.6em}

  \begin{columns}[T,onlytextwidth]
    \begin{column}{0.47\textwidth}
      \begin{block}{Chatbot}
        Guesses a number from what it remembers about the 1994 paper.
      \end{block}
    \end{column}
    \begin{column}{0.47\textwidth}
      \begin{block}{Agent}
        Runs \artifact{did\_analysis.py} on our folder and quotes the printed note.
      \end{block}
    \end{column}
  \end{columns}
  \vspace{0.6em}

  Only one of these answers comes from \textbf{your} data.
\end{frame}

\begin{frame}{C1 \textperiodcentered\ Your turn}
  \youtry{read only, nothing is edited}
  \vspace{0.3em}
  \begin{enumerate}
    \item Ask a plain chat: ``What is the NJ--PA change in employment?''
    \item Ask the agent the same question, adding ``Do not edit any file.''
    \item Compare: where did each number come from?
  \end{enumerate}
  \vspace{0.5em}
  {\small Guide: \artifact{demos/slot1\_workflows.md} \textperiodcentered\ \artifact{exercises/04\_slot1\_workflows.md}}
\end{frame}

% ---------------------------------------------------------------- C2
\begin{frame}{C2 \textperiodcentered\ Version control}
  \timebox{09:30--09:55}
  \inplainwords{A record of every saved version of the project, each with a note,
    so you can see what changed and go back. GitHub keeps that record online and
    lets others review changes before they are accepted.}
  \likethis{Track Changes plus dated ``save as'' copies, shared with coauthors.}
\end{frame}

\begin{frame}{C2 \textperiodcentered\ Five words we use all day}
  \begin{tabular}{@{}p{0.27\linewidth}p{0.66\linewidth}@{}}
    \toprule
    \textbf{Word} & \textbf{Meaning} \\
    \midrule
    Project folder & All the course files, with their full history \\
    Saved version & A snapshot of the files, with a one-line note \\
    Separate copy & A safe place to try changes (a ``branch'') \\
    List of changes & The lines added and removed, side by side \\
    Pull request & A request to merge changes, with a review page \\
    \bottomrule
  \end{tabular}
\end{frame}

\begin{frame}{C2 \textperiodcentered\ Card--Krueger: reading the agent's edit}
  \begin{itemize}
    \item The agent edits \artifact{examples/card-krueger-toy/README.md}.
    \item GitHub shows a list of changes: green lines added, red lines removed.
    \item You accept, or ask for a fix, before it joins the main version.
    \item Nothing is lost: every saved version can come back.
  \end{itemize}
\end{frame}

\begin{frame}{C2 \textperiodcentered\ Your turn: first pull request}
  \youtry{09:30--09:55}
  \vspace{0.3em}
  \begin{enumerate}
    \item Make a separate copy.
    \item Edit one line in \artifact{notes/lane\_choice.md}: your tool and why.
    \item Save the version with a short note.
    \item Open a pull request and read its list of changes.
    \item Review, then merge.
  \end{enumerate}
  {\small Guide: \artifact{sessions/slot1.md} \textperiodcentered\ \artifact{exercises/05\_slot1\_governance.md}}
\end{frame}

% ---------------------------------------------------------------- C3
\begin{frame}{C3 \textperiodcentered\ \texttt{AGENTS.md}: boundaries first}
  \timebox{09:55--10:35}
  \inplainwords{A plain-text rules file the agent reads before it starts any task.
    It says what the agent may touch, what it must never do, and how to check its work.}
  \likethis{the house rules you give a new RA on day one.}
\end{frame}

\begin{frame}{C3 \textperiodcentered\ Card--Krueger: the agent refuses}
  Two lines from our \artifact{AGENTS.md}:
  \vspace{0.3em}
  \begin{block}{}
    \texttt{- Do not modify raw data.}\\
    \texttt{- Do not let agents certify identification, causality, or interpretation.}
  \end{block}
  \vspace{0.4em}
  Demo: ask the agent to edit \artifact{toy\_fast\_food.csv}.\\
  Watch it refuse and point to the rule.
\end{frame}

\begin{frame}{C3 \textperiodcentered\ Privacy in three questions}
  \begin{enumerate}
    \item How sensitive is the data?
    \item Who may see it?
    \item Where does the agent run: your laptop or someone else's computer?
  \end{enumerate}
  \vspace{0.5em}
  \begin{block}{Faculty default}
    \centering\large Anonymize \textperiodcentered\ authorize \textperiodcentered\ keep it on the machine
  \end{block}
  \vspace{0.3em}
  {\small Full check: \artifact{docs/privacy\_boundary.md}}
\end{frame}

\begin{frame}{C3 \textperiodcentered\ Your turn}
  \youtry{09:55--10:35}
  \vspace{0.3em}
  \begin{enumerate}
    \item Read \artifact{AGENTS.md} (two minutes).
    \item Add one rule you would want for your own project.
    \item Fill in the privacy check in \artifact{docs/privacy\_boundary.md}.
  \end{enumerate}
  \vspace{0.5em}
  {\small Guide: \artifact{demos/slot1\_governance.md} \textperiodcentered\ Pause 10:35--10:45}
\end{frame}

% ---------------------------------------------------------------- C4
\begin{frame}{C4 \textperiodcentered\ To-do lists}
  \timebox{10:45--11:00}
  \inplainwords{Before acting, the agent writes the steps it plans to take as a
    checklist. You edit the list, then it works through it and ticks items off.}
  \likethis{a shopping list you check before anyone drives to the store.}
\end{frame}

\begin{frame}{C4 \textperiodcentered\ Card--Krueger: plan only}
  \prompttext{\raggedright Plan only, do not edit. Task: examples/card-krueger-goals/goals/ck-robustness-checks.md}
  \vspace{0.6em}
  \begin{itemize}
    \item[$\square$] Load the panel
    \item[$\square$] Check that every store appears before and after
    \item[$\square$] Compute NJ and PA group means
    \item[$\square$] Drop one store at a time and recompute
    \item[$\square$] Write a short note
  \end{itemize}
\end{frame}

\begin{frame}{C4 \textperiodcentered\ Your turn}
  \youtry{10:45--11:00}
  \vspace{0.3em}
  \begin{enumerate}
    \item Ask for the to-do list only.
    \item Delete or reorder one step.
    \item Let the agent do step 1 only, then stop.
    \item Paste the list in \artifact{notes/modes\_playbook.md}.
  \end{enumerate}
  \vspace{0.5em}
  {\small Template: \artifact{templates/agent\_task\_contract.md}}
\end{frame}

% Slot 2: C5--C8 (Thu 26 Nov 11:00--12:00 + Fri 27 Nov 09:30--10:15 Rome)

\section{Slot 2: C5--C8}

\begin{frame}{Slot 2 \textperiodcentered\ Thu 11:00--12:00 + Fri 09:30--10:15}
  \begin{center}
    \begin{tabular}{@{}rl@{}}
      \toprule
      \textbf{Clock} & \textbf{Block} \\
      \midrule
      Thu 11:00--11:25 & \concept{C5} Skills: why \\
      11:25--11:55 & \concept{C6} Create a Skill \\
      11:55--12:00 & Thursday wrap, overnight task \\
      \addlinespace[0.3em]
      Fri 09:30--09:35 & Recap \\
      09:35--09:55 & \concept{C7} Sub-agent vs Skill \\
      09:55--10:15 & \concept{C8} Rule of thumb: Skill or sub-agent? \\
      \bottomrule
    \end{tabular}
  \end{center}
\end{frame}

% ---------------------------------------------------------------- C5
\begin{frame}{C5 \textperiodcentered\ Skills: why}
  \timebox{Thu 11:00--11:25}
  \inplainwords{A Skill is a saved set of instructions for a task you do often.
    The agent follows it the same way every time, and you can share it with coauthors.}
  \likethis{a recipe card, in place of explaining the dish from memory each time.}
\end{frame}

\begin{frame}{C5 \textperiodcentered\ Card--Krueger: \texttt{replication-checker}}
  Does our folder reproduce its numbers from a fresh start?
  \vspace{0.4em}
  \begin{itemize}
    \item File: \artifact{.claude/skills/replication-checker/SKILL.md}
    \item Inside: a name, when to use it, numbered steps.
    \item It runs the check command and confirms the data file is unchanged.
    \item It reports \textcolor{codegreen}{\textbf{GREEN}}, \textbf{YELLOW} or \textcolor{codered}{\textbf{RED}}, with the output as evidence.
    \item It never certifies causality.
  \end{itemize}
\end{frame}

\begin{frame}{C5 \textperiodcentered\ Your turn}
  \youtry{Thu 11:00--11:25}
  \vspace{0.3em}
  \begin{enumerate}
    \item Read the Skill file before you trust it: what may it run?
    \item Run it on \artifact{examples/card-krueger-toy/}.
    \item Write down its colour and the one line of evidence behind it.
  \end{enumerate}
  \vspace{0.5em}
  {\small Guide: \artifact{demos/slot2\_sdd\_skill\_pr.md} \textperiodcentered\ \artifact{exercises/06\_slot2\_sdd\_skill\_pr.md}}
\end{frame}

% ---------------------------------------------------------------- C6
\begin{frame}{C6 \textperiodcentered\ Create a Skill}
  \timebox{Thu 11:25--11:55}
  \inplainwords{Make a folder and write a short file in it with a name, a line saying
    when to use it, and numbered steps. Test it once and fix whatever was unclear.}
  \likethis{writing down your referee-report checklist so a student can follow it.}
\end{frame}

\begin{frame}[fragile]{C6 \textperiodcentered\ Card--Krueger: \texttt{sample-restrictions-note}}
  {\small \artifact{.claude/skills/sample-restrictions-note/SKILL.md} (shape of the file)}
  \begin{lstlisting}[style=markdown,basicstyle=\ttfamily\footnotesize]
---
name: sample-restrictions-note
description: Use after any change to the Card-Krueger sample.
---
1. List every rule that drops stores.
2. Give the store count after each step.
3. Quote the command that printed the counts (panel_workflow.py).
4. Never edit the data file.
  \end{lstlisting}
  {\small The header lines between the dashes: its name and when to use it.}
\end{frame}

\begin{frame}{C6 \textperiodcentered\ Your turn}
  \youtry{Thu 11:25--11:55}
  \vspace{0.3em}
  \begin{enumerate}
    \item Fill in \artifact{templates/skill\_template.md}.
    \item On your separate copy, save it as \artifact{.claude/skills/<name>/SKILL.md}.
    \item Run it once on the Card--Krueger folder.
    \item Fix the one step the agent misread.
  \end{enumerate}
\end{frame}

\begin{frame}{Thursday wrap and Friday recap}
  \timebox{Thu 11:55--12:00 \textperiodcentered\ Fri 09:30--09:35}
  \vspace{0.3em}
  \begin{block}{Tonight}
    Draft one Skill for a task you repeat every week (\artifact{docs/between\_sessions.md}).
  \end{block}
  \vspace{0.4em}
  \begin{block}{Friday 09:30: so far}
    \concept{C1} agent \textperiodcentered\ \concept{C2} saved versions \textperiodcentered\
    \concept{C3} rules file \textperiodcentered\ \concept{C4} to-do list \textperiodcentered\
    \concept{C5--C6} Skills
  \end{block}
\end{frame}

% ---------------------------------------------------------------- C7
\begin{frame}{C7 \textperiodcentered\ Sub-agent: what it is}
  \timebox{Fri 09:35--09:55}
  \inplainwords{A sub-agent is a separate helper agent with its own blank notebook and
    one narrow job. It reports back a summary. A Skill is instructions; a sub-agent
    is another worker.}
  \likethis{a recipe (Skill) vs a sous-chef you send to taste the sauce (sub-agent).}
\end{frame}

\begin{frame}{C7 \textperiodcentered\ Card--Krueger: Skill vs sub-agent}
  \begin{tabular}{@{}p{0.22\linewidth}p{0.34\linewidth}p{0.36\linewidth}@{}}
    \toprule
     & \textbf{Skill} & \textbf{Sub-agent} \\
    \midrule
    What it is & Instructions & Another worker \\
    Memory & Shares your chat & Its own blank notebook \\
    Gives back & The finished task & A short report \\
    Our example & \artifact{replication-checker} & \artifact{data-reviewer} \\
    \bottomrule
  \end{tabular}
  \vspace{0.6em}

  \artifact{data-reviewer} looks only at panel balance and NJ/PA coding. It never edits.
\end{frame}

\begin{frame}{C7 \textperiodcentered\ Your turn}
  \youtry{Fri 09:35--09:55}
  \vspace{0.3em}
  Open \artifact{.claude/agents/data-reviewer.md} and mark three things:
  \begin{enumerate}
    \item \textbf{The job:} what does it look at?
    \item \textbf{The limits:} what is it not allowed to do?
    \item \textbf{The report:} what does it hand back?
  \end{enumerate}
\end{frame}

% ---------------------------------------------------------------- C8
\begin{frame}{C8 \textperiodcentered\ Rule of thumb: Skill or sub-agent?}
  \timebox{Fri 09:55--10:15}
  \inplainwords{Same steps done the same way: a \textbf{Skill}.\\
    Fresh pair of eyes, a big messy job, or work that can run at the same time:
    a \textbf{sub-agent}.}
  \likethis{a checklist for yourself vs asking a colleague to referee.}
\end{frame}

\begin{frame}{C8 \textperiodcentered\ Card--Krueger: two tasks}
  \begin{columns}[T,onlytextwidth]
    \begin{column}{0.47\textwidth}
      \begin{block}{Skill}
        ``Recompute the DiD the same way each time.''
      \end{block}
    \end{column}
    \begin{column}{0.47\textwidth}
      \begin{block}{Sub-agent}
        ``Independently check my data coding.''
      \end{block}
    \end{column}
  \end{columns}
  \vspace{0.8em}
  You stay in the main chat: you decide, approve and merge.
\end{frame}

% Answer key: 1 Skill, 2 sub-agent, 3 Skill, 4 sub-agent, 5 Skill, 6 sub-agent.
\begin{frame}{C8 \textperiodcentered\ Your turn: sort six tasks}
  \youtry{Fri 09:55--10:15 \textperiodcentered\ then compare with a neighbour}
  \vspace{0.2em}
  \small
  \begin{tabular}{@{}rp{0.68\linewidth}l@{}}
    \toprule
    & \textbf{Card--Krueger task} & \textbf{Skill or sub-agent?} \\
    \midrule
    1 & Recompute group means after every data change & \\
    2 & Check NJ/PA coding without seeing my code & \\
    3 & Write the sample-restrictions note the same way & \\
    4 & Review literature claims while I fix code & \\
    5 & Check the folder reproduces from a fresh start & \\
    6 & Referee my pull request with fresh eyes & \\
    \bottomrule
  \end{tabular}
\end{frame}

% Slot 3: C9--C11 launch (Fri 27 Nov 10:15--12:00 Rome)

\section{Slot 3: C9--C11}

\begin{frame}{Slot 3 \textperiodcentered\ Fri 10:15--12:00}
  \begin{center}
    \begin{tabular}{@{}rl@{}}
      \toprule
      \textbf{Clock} & \textbf{Block} \\
      \midrule
      10:15--10:50 & \concept{C9} Using sub-agents \\
      10:50--11:15 & \concept{C10} Orchestrating sub-agents and swarms \\
      11:15--11:25 & Pause \\
      11:25--12:00 & \concept{C11} Issue-board swarm: launch two cards \\
      \textit{12:00--14:00} & \textit{Lunch: the agents keep working} \\
      \bottomrule
    \end{tabular}
  \end{center}
\end{frame}

% ---------------------------------------------------------------- C9
\begin{frame}{C9 \textperiodcentered\ Using sub-agents}
  \timebox{10:15--10:50}
  \inplainwords{Call the helper by name, give it one narrow instruction, keep it
    read-only and read its report before you act. You decide what to do with the report.}
  \likethis{sending a draft to two referees who don't talk to each other.}
\end{frame}

\begin{frame}{C9 \textperiodcentered\ Card--Krueger: two referee reports}
  \prompttext{\raggedright Use the data-reviewer sub-agent on examples/card-krueger-toy/. Read only. Report panel balance and NJ/PA coding.}
  \vspace{0.4em}
  \prompttext{\raggedright Use the literature-reviewer sub-agent on the same folder. Read only. Check the claims against the references.}
  \vspace{0.6em}
  \begin{itemize}
    \item Two separate notebooks: neither sees the other's report.
    \item Where they agree, look closer. Where they disagree, you decide.
  \end{itemize}
\end{frame}

\begin{frame}{C9 \textperiodcentered\ Your turn}
  \youtry{10:15--10:50}
  \vspace{0.3em}
  \begin{enumerate}
    \item Run both helpers on the Card--Krueger folder.
    \item Log one finding you accept and one you reject in \artifact{notes/orchestration\_log.md}.
    \item Check that no file changed.
  \end{enumerate}
  \vspace{0.4em}
  {\small Helper won't start? Paste its instructions into a fresh chat and say ``read only''.}\\
  {\small Guide: \artifact{demos/slot3\_subagents\_vc.md} \textperiodcentered\ \artifact{exercises/07\_slot3\_subagents\_vc.md}}
\end{frame}

% ---------------------------------------------------------------- C10
\begin{frame}{C10 \textperiodcentered\ Orchestrating sub-agents and swarms}
  \timebox{10:50--11:15}
  \inplainwords{Orchestrating means splitting a job into parts, handing each part to a
    different agent, and putting the results back together. A swarm is many agents
    working at the same time; a human still makes the final call.}
  \likethis{a PI running several RAs, each with their own files.}
\end{frame}

\begin{frame}{C10 \textperiodcentered\ Card--Krueger: a writer and a checker}
  \begin{center}
    \begin{tikzpicture}[
        node distance=1.0cm,
        role/.style={rectangle, draw=courseblue, thick, rounded corners=3pt, minimum width=3.3cm, minimum height=1.0cm, align=center, font=\footnotesize\bfseries, fill=coursebg},
        arr/.style={-Stealth, thick, courseblue}
      ]
      \node[role] (w) {Writer\\\normalfont\scriptsize documents the panel};
      \node[role, right=of w, fill=courseteal!14] (c) {Checker\\\normalfont\scriptsize reviews literature claims};
      \node[role, below=0.9cm of $(w)!0.5!(c)$, draw=codered, fill=codered!8] (y) {You\\\normalfont\scriptsize read both, then merge};
      \draw[arr] (w) -- (y);
      \draw[arr] (c) -- (y);
    \end{tikzpicture}
  \end{center}
  \vspace{0.3em}
  \begin{itemize}
    \item Different files: the two agents never edit the same lines.
    \item Rules for the pair: \artifact{examples/card-krueger-swarm/README.md}.
  \end{itemize}
\end{frame}

\begin{frame}{C10 \textperiodcentered\ Your turn: split the job}
  \youtry{10:50--11:15}
  \vspace{0.3em}
  Split ``replicate and extend Card--Krueger'' into three jobs that share no files.
  \vspace{0.5em}

  \begin{tabular}{@{}lp{0.40\linewidth}p{0.32\linewidth}@{}}
    \toprule
    \textbf{Job} & \textbf{Files it may touch} & \textbf{Done when} \\
    \midrule
    1 & & \\
    2 & & \\
    3 & & \\
    \bottomrule
  \end{tabular}
  \vspace{0.5em}

  {\small Pause 11:15--11:25.}
\end{frame}

% ---------------------------------------------------------------- C11
\begin{frame}{C11 \textperiodcentered\ Issue-board swarm}
  \timebox{11:25--12:00}
  \inplainwords{Each job is written as an issue, a to-do card on GitHub. Agents pick
    up cards marked ready, work in their own copy, and send back a pull request
    linked to the card for review.}
  \likethis{a job board in a shared office: whoever is free takes the next card.}
\end{frame}

\begin{frame}{C11 \textperiodcentered\ Card--Krueger: two cards on the board}
  \begin{center}
    \begin{tikzpicture}[
        node distance=0.55cm,
        st/.style={rectangle, draw=courseblue, thick, rounded corners=3pt, minimum width=2.4cm, minimum height=0.8cm, align=center, font=\footnotesize\bfseries, fill=coursebg},
        arr/.style={-Stealth, thick, courseblue}
      ]
      \node[st] (r) {ready};
      \node[st, right=of r] (p) {in-progress};
      \node[st, right=of p, fill=courseteal!14] (n) {needs-review};
      \node[st, right=of n, fill=codegreen!12] (d) {done};
      \draw[arr] (r) -- (p);
      \draw[arr] (p) -- (n);
      \draw[arr] (n) -- (d);
    \end{tikzpicture}
  \end{center}
  \vspace{0.4em}
  \begin{itemize}
    \item Card 1: \artifact{ck-data-cleaning}: check and document the panel.
    \item Card 2: \artifact{ck-parallel-trends-figure}: draw the before/after figure.
    \item Cards live in \artifact{examples/card-krueger-swarm/issues/}.
    \item Only you move a card to \textbf{done}.
  \end{itemize}
\end{frame}

\begin{frame}{C11 \textperiodcentered\ Your turn: launch before lunch}
  \youtry{11:25--12:00}
  \vspace{0.3em}
  \begin{enumerate}
    \item Cloud agent? First fill in \artifact{agent-harness/autonomous\_agent\_risk\_card.md}.
    \item Open the two cards as issues and mark them \textbf{ready}.
    \item Start two agents, one per card (cloud, or two local windows).
    \item Go to lunch. Back at 14:00 to review.
  \end{enumerate}
  \vspace{0.4em}
  {\small At most two cards live. No confidential files. Agents never merge.}
\end{frame}

% Slot 4: Capstone, close (Fri 27 Nov 14:00--16:30 Rome)

\section{Slot 4}

\begin{frame}{Slot 4 \textperiodcentered\ Fri 14:00--16:30}
  \begin{center}
    \begin{tabular}{@{}rl@{}}
      \toprule
      \textbf{Clock} & \textbf{Block} \\
      \midrule
      14:00--14:10 & Review what the swarm returned \\
      14:10--14:25 & Goals \\
      14:25--14:50 & Acceptance criteria: rewrite one card, relaunch \\
      14:50--15:50 & Verification and evaluation \\
      15:50--16:15 & Capstone: your own project \\
      16:15--16:30 & Close and 7-day plan \\
      \bottomrule
    \end{tabular}
  \end{center}
\end{frame}

% ---------------------------------------------------------------- review
\begin{frame}{Review what came back}
  \timebox{14:00--14:10}\par\vspace{0.3em}
  For each pull request from lunch:
  \begin{enumerate}
    \item Ask \artifact{pr-reviewer} (read only) for a report.
    \item Did it touch only the files on its card?
    \item Does the check command pass?
    \item Move the card to \textbf{done}, or back to \textbf{ready} with a note.
  \end{enumerate}
  \vspace{0.4em}
  {\large How did each agent know when to stop?}
\end{frame}

% ----------------------------------------------------------------
\begin{frame}{Goals}
  \timebox{14:10--14:25}
  \inplainwords{A goal is one sentence describing the finished state, in a form you
    can check. It tells the agent where to stop, not how to get there.}
  \likethis{the destination in a satnav, not the turn-by-turn directions.}
\end{frame}

\begin{frame}{Card--Krueger: from vague to checkable}
  \begin{block}{Vague}
    ``Audit the replication.''
  \end{block}
  \vspace{0.3em}
  \begin{block}{Goal}
    \artifact{notes/replication\_audit\_ck.md} exists and records GREEN, YELLOW or RED
    with the check output.
  \end{block}
  \vspace{0.4em}
  {\small Claude Code and Codex: type \texttt{/goal} first. Cursor: make it the last line of the prompt.}\\
  {\small More goals: \artifact{examples/card-krueger-goals/goals/}}
\end{frame}

\begin{frame}{Your turn}
  \youtry{14:10--14:25}
  \vspace{0.3em}
  \begin{enumerate}
    \item Take one lunchtime card.
    \item Rewrite its title as a one-sentence goal.
    \item Test: could a stranger answer ``done?'' with yes or no?
  \end{enumerate}
\end{frame}

% ----------------------------------------------------------------
\begin{frame}{Acceptance criteria}
  \timebox{14:25--14:50}
  \inplainwords{A yes/no check that says when a job counts as done. In a swarm, each
    agent knows when to stop, and you can check many results quickly.}
  \likethis{a building inspector's checklist.}
\end{frame}

\begin{frame}{Card--Krueger: done when\ldots}
  Card \artifact{ck-data-cleaning} is done when:
  \vspace{0.3em}
  \begin{itemize}
    \item[$\square$] \artifact{toy\_fast\_food.csv} is unchanged.
    \item[$\square$] The check command passes.
    \item[$\square$] The README states the sample restriction.
    \item[$\square$] The DiD printed equals the DiD in the note.
  \end{itemize}
  \vspace{0.4em}
  Each line is yes or no. No line asks ``is it good?''
\end{frame}

\begin{frame}{Your turn: rewrite and relaunch}
  \youtry{14:25--14:50}
  \vspace{0.3em}
  \begin{enumerate}
    \item Add three yes/no criteria to one card.
    \item Relaunch the agent on that card.
    \item Judge the new pull request only against your three lines.
  \end{enumerate}
  \vspace{0.4em}
  {\small Template: \artifact{templates/issue\_card.md}}
\end{frame}

% ----------------------------------------------------------------
\begin{frame}{Verification and evaluation}
  \timebox{14:50--15:50}
  \begin{enumerate}
    \item \textbf{Check less:} small jobs, protected files, clear done-when checks.
    \item \textbf{Automate:} check commands that run on their own, on every change.
    \item \textbf{Test your set-up:} try your Skills, sub-agents and whole set-up
      on known mistakes.
  \end{enumerate}
  \likethis{planting a known error in a paper to see whether the referee catches it.}
\end{frame}

\begin{frame}{Card--Krueger: a planted mistake}
  \artifact{faults/drifted\_table.tex}: a hand-copied DiD (3.15) and row count (47).
  \vspace{0.5em}
  \begin{center}
    \begin{tikzpicture}[
        node distance=0.7cm,
        vbox/.style={rectangle, draw=courseblue, thick, rounded corners=3pt, minimum width=2.5cm, minimum height=0.85cm, align=center, font=\footnotesize\bfseries, fill=coursebg},
        arr/.style={-Stealth, thick, courseblue}
      ]
      \node[vbox, draw=codered, fill=codered!8] (f) {1. Fail\\\normalfont\scriptsize check stops};
      \node[vbox, right=of f] (d) {2. Diagnose\\\normalfont\scriptsize which number?};
      \node[vbox, right=of d] (x) {3. Fix\\\normalfont\scriptsize never the checks};
      \node[vbox, right=of x, fill=codegreen!12] (r) {4. Rerun\\\normalfont\scriptsize check passes};
      \draw[arr] (f) -- (d);
      \draw[arr] (d) -- (x);
      \draw[arr] (x) -- (r);
    \end{tikzpicture}
  \end{center}
  \vspace{0.4em}
  A passing check means the numbers match. It says nothing about causality.
\end{frame}

\begin{frame}{Card--Krueger: test the helpers}
  \begin{tabular}{@{}p{0.29\linewidth}p{0.43\linewidth}p{0.20\linewidth}@{}}
    \toprule
    \textbf{What you test} & \textbf{Known mistake} & \textbf{Result} \\
    \midrule
    Skill:\newline \artifact{replication-checker} & \artifact{drifted\_table.tex} & caught / missed \\
    Sub-agent:\newline \artifact{data-reviewer} & \artifact{unbalanced\_panel.csv} & caught / missed \\
    Whole set-up & Fresh chat, fresh copy: reproduce from the README alone? & yes / no \\
    \bottomrule
  \end{tabular}
  \vspace{0.6em}

  {\small Planted mistakes: \artifact{examples/card-krueger-toy/faults/}}
\end{frame}

\begin{frame}{Your turn}
  \youtry{14:50--15:50}
  \vspace{0.3em}
  \begin{enumerate}
    \item On a temporary separate copy, fix one planted mistake.
    \item Do not touch \artifact{tests/}. Rerun the check command.
    \item Fill in the three-row scorecard: Skill, sub-agent, whole set-up.
    \item Record it in \artifact{templates/agent\_use\_record.md}.
  \end{enumerate}
  \vspace{0.3em}
  {\small Optional demo: an automatic check after each edit (\artifact{.claude/hooks/}).}\\
  {\small Guide: \artifact{demos/slot4\_verify\_capstone.md}}
\end{frame}

% ---------------------------------------------------------------- Capstone and close
\section{Capstone and close}

\begin{frame}{Capstone: your own project}
  \timebox{15:50--16:15}\par\vspace{0.3em}
  On one non-confidential project of yours:
  \begin{enumerate}
    \item One rule in its \artifact{AGENTS.md}.
    \item One Skill for a task you repeat.
    \item One goal with three yes/no criteria.
    \item Evidence: the check command and its output.
  \end{enumerate}
  {\small Only confidential data at hand? Use the Card--Krueger folder. Rubric: \artifact{docs/capstone\_rubric.md}}
\end{frame}

\begin{frame}{Close: your next seven days}
  \timebox{16:15--16:30}\par\vspace{0.3em}
  \begin{itemize}
    \item Day 1: add \artifact{AGENTS.md} to one real project.
    \item This week: turn one weekly chore into a Skill.
    \item Next job you delegate: write a goal and three done-when checks first.
    \item Always: read the list of changes; you sign off, not the agent.
  \end{itemize}
  \vspace{0.4em}
  {\small Plan: \artifact{docs/adoption\_plan\_7d.md} \textperiodcentered\ Start page: \artifact{START\_HERE.md}}
\end{frame}


\end{document}
